\IfFileExists{stacks-project.cls}{%
\documentclass{stacks-project}
}{%
\documentclass{amsart}
}
\usepackage{amssymb}

% For dealing with references we use the comment environment
\usepackage{verbatim}
\newenvironment{reference}{\comment}{\endcomment}
%\newenvironment{reference}{}{}
\newenvironment{slogan}{\comment}{\endcomment}
\newenvironment{history}{\comment}{\endcomment}

% For commutative diagrams we use Xy-pic
\usepackage[all]{xy}

% We use 2cell for 2-commutative diagrams.
\xyoption{2cell}
\UseAllTwocells

% We use multicol for the list of chapters between chapters
\usepackage{multicol}

% This is generally recommended for better output
\usepackage{lmodern}
\usepackage[T1]{fontenc}
\usepackage{textalpha}

% For cross-file-references
\usepackage{xr-hyper}
\newcommand{\maybeexternaldocument}[2]{%
  \IfFileExists{#2.aux}{%
    \externaldocument[#1]{#2}%
  }{%
    \PackageWarning{stack}{Missing #2.aux; external references with prefix #1 unavailable}%
  }%
}
\let\externaldocumentorig\externaldocument
\renewcommand{\externaldocument}[2][]{%
  \IfFileExists{#2.aux}{%
    \externaldocumentorig[#1]{#2}%
  }{%
    \PackageWarning{stack}{Missing #2.aux; external references with prefix #1 unavailable}%
  }%
}

% Package for hypertext links:
\usepackage{hyperref}

% For any local file, say "hello.tex" you want to link to please
% use \externaldocument[hello-]{hello}
\externaldocument[introduction-]{introduction}
\externaldocument[conventions-]{conventions}
\externaldocument[sets-]{sets}
\externaldocument[categories-]{categories}
\externaldocument[topology-]{topology}
\externaldocument[sheaves-]{sheaves}
\externaldocument[sites-]{sites}
\externaldocument[stacks-]{stacks}
\externaldocument[fields-]{fields}
\externaldocument[algebra-]{algebra}
\externaldocument[brauer-]{brauer}
\externaldocument[homology-]{homology}
\externaldocument[derived-]{derived}
\externaldocument[simplicial-]{simplicial}
\externaldocument[more-algebra-]{more-algebra}
\externaldocument[smoothing-]{smoothing}
\externaldocument[modules-]{modules}
\externaldocument[sites-modules-]{sites-modules}
\externaldocument[injectives-]{injectives}
\externaldocument[cohomology-]{cohomology}
\externaldocument[sites-cohomology-]{sites-cohomology}
\externaldocument[dga-]{dga}
\externaldocument[dpa-]{dpa}
\externaldocument[sdga-]{sdga}
\externaldocument[hypercovering-]{hypercovering}
\externaldocument[schemes-]{schemes}
\externaldocument[constructions-]{constructions}
\externaldocument[properties-]{properties}
\externaldocument[morphisms-]{morphisms}
\externaldocument[coherent-]{coherent}
\externaldocument[divisors-]{divisors}
\externaldocument[limits-]{limits}
\externaldocument[varieties-]{varieties}
\externaldocument[topologies-]{topologies}
\externaldocument[descent-]{descent}
\externaldocument[perfect-]{perfect}
\externaldocument[more-morphisms-]{more-morphisms}
\externaldocument[flat-]{flat}
\externaldocument[groupoids-]{groupoids}
\externaldocument[more-groupoids-]{more-groupoids}
\externaldocument[etale-]{etale}
\externaldocument[chow-]{chow}
\externaldocument[intersection-]{intersection}
\externaldocument[pic-]{pic}
\externaldocument[weil-]{weil}
\externaldocument[adequate-]{adequate}
\externaldocument[dualizing-]{dualizing}
\externaldocument[duality-]{duality}
\externaldocument[discriminant-]{discriminant}
\externaldocument[derham-]{derham}
\externaldocument[local-cohomology-]{local-cohomology}
\externaldocument[algebraization-]{algebraization}
\externaldocument[curves-]{curves}
\externaldocument[resolve-]{resolve}
\externaldocument[models-]{models}
\externaldocument[functors-]{functors}
\externaldocument[equiv-]{equiv}
\externaldocument[pione-]{pione}
\externaldocument[etale-cohomology-]{etale-cohomology}
\externaldocument[proetale-]{proetale}
\externaldocument[relative-cycles-]{relative-cycles}
\externaldocument[more-etale-]{more-etale}
\externaldocument[trace-]{trace}
\externaldocument[crystalline-]{crystalline}
\externaldocument[spaces-]{spaces}
\externaldocument[spaces-properties-]{spaces-properties}
\externaldocument[spaces-morphisms-]{spaces-morphisms}
\externaldocument[decent-spaces-]{decent-spaces}
\externaldocument[spaces-cohomology-]{spaces-cohomology}
\externaldocument[spaces-limits-]{spaces-limits}
\externaldocument[spaces-divisors-]{spaces-divisors}
\externaldocument[spaces-over-fields-]{spaces-over-fields}
\externaldocument[spaces-topologies-]{spaces-topologies}
\externaldocument[spaces-descent-]{spaces-descent}
\externaldocument[spaces-perfect-]{spaces-perfect}
\externaldocument[spaces-more-morphisms-]{spaces-more-morphisms}
\externaldocument[spaces-flat-]{spaces-flat}
\externaldocument[spaces-groupoids-]{spaces-groupoids}
\externaldocument[spaces-more-groupoids-]{spaces-more-groupoids}
\externaldocument[bootstrap-]{bootstrap}
\externaldocument[spaces-pushouts-]{spaces-pushouts}
\externaldocument[spaces-chow-]{spaces-chow}
\externaldocument[groupoids-quotients-]{groupoids-quotients}
\externaldocument[spaces-more-cohomology-]{spaces-more-cohomology}
\externaldocument[spaces-simplicial-]{spaces-simplicial}
\externaldocument[spaces-duality-]{spaces-duality}
\externaldocument[formal-spaces-]{formal-spaces}
\externaldocument[restricted-]{restricted}
\externaldocument[spaces-resolve-]{spaces-resolve}
\externaldocument[formal-defos-]{formal-defos}
\externaldocument[defos-]{defos}
\externaldocument[cotangent-]{cotangent}
\externaldocument[examples-defos-]{examples-defos}
\externaldocument[algebraic-]{algebraic}
\externaldocument[examples-stacks-]{examples-stacks}
\externaldocument[stacks-sheaves-]{stacks-sheaves}
\externaldocument[criteria-]{criteria}
\externaldocument[artin-]{artin}
\externaldocument[quot-]{quot}
\externaldocument[stacks-properties-]{stacks-properties}
\externaldocument[stacks-morphisms-]{stacks-morphisms}
\externaldocument[stacks-limits-]{stacks-limits}
\externaldocument[stacks-cohomology-]{stacks-cohomology}
\externaldocument[stacks-perfect-]{stacks-perfect}
\externaldocument[stacks-introduction-]{stacks-introduction}
\externaldocument[stacks-more-morphisms-]{stacks-more-morphisms}
\externaldocument[stacks-geometry-]{stacks-geometry}
\externaldocument[moduli-]{moduli}
\externaldocument[moduli-curves-]{moduli-curves}
\externaldocument[examples-]{examples}
\externaldocument[exercises-]{exercises}
\externaldocument[guide-]{guide}
\externaldocument[desirables-]{desirables}
\externaldocument[coding-]{coding}
\externaldocument[obsolete-]{obsolete}
\externaldocument[fdl-]{fdl}
\externaldocument[index-]{index}

% Heap Project documents
\externaldocument[linear-algebra-]{linear-algebra}
\externaldocument[groups-]{groups}
\externaldocument[complex-analysis-]{complex-analysis}
\externaldocument[functional-analysis-]{functional-analysis}
\externaldocument[categories-]{categories}
\externaldocument[commutative-algebra-]{commutative-algebra}
\externaldocument[ringed-spaces-]{ringed-spaces}
\externaldocument[homological-algebra-]{homological-algebra}
\externaldocument[lie-algebras-]{lie-algebras}
\externaldocument[representation-theory-]{representation-theory}
\externaldocument[smooth-manifolds-]{smooth-manifolds}
\externaldocument[differential-topology-]{differential-topology}
\externaldocument[algebraic-topology-]{algebraic-topology}
\externaldocument[differential-geometry-]{differential-geometry}
\externaldocument[algebraic-geometry-]{algebraic-geometry}
\externaldocument[symplectic-geometry-]{symplectic-geometry}
\externaldocument[complex-geometry-]{complex-geometry}
\externaldocument[hodge-theory-]{hodge-theory}
\maybeexternaldocument{deformations-}{deformations}
\externaldocument[vertex-algebras-]{vertex-algebras}
\externaldocument[springer-]{springer}
\externaldocument[infty-categories-]{infty-categories}
\externaldocument[seminars-differential-geometry-]{%
seminars-differential-geometry}
\externaldocument[seminars-complex-geometry-]{%
seminars-complex-geometry}
\externaldocument[seminars-algebraic-geometry-]{%
seminars-algebraic-geometry}
\externaldocument[seminars-others-]{%
seminars-others}
%\externaldocument[geometric-prequantization-]{geometric-prequantization}
\externaldocument[surfaces-]{surfaces}
\externaldocument[birational-maps-commutative-algebra-]{birational-maps-commutative-algebra}
\externaldocument[cremona-transformations-]{cremona-transformations}
\externaldocument[derived-categories-of-sheaves-]{derived-categories-of-sheaves}
\externaldocument[geometric-invariant-theory-]{geometric-invariant-theory}
\externaldocument[geometric-stability-conditions-and-group-actions-]{geometric-stability-conditions-and-group-actions}
\externaldocument[moduli-spaces-of-sheaves-]{moduli-spaces-of-sheaves}
\externaldocument[bridgeland-stability-general-theory-]{bridgeland-stability-general-theory}
\externaldocument[hilbert-schemes-]{hilbert-schemes}
\externaldocument[noncommutative-abelian-surfaces-and-kummer-type-hyperkahler-varieties-]{noncommutative-abelian-surfaces-and-kummer-type-hyperkahler-varieties}
\externaldocument[mumford-tate-groups-in-hodge-theory-]{mumford-tate-groups-in-hodge-theory}
\externaldocument[hodge-birational-atoms-2026-]{hodge-birational-atoms-2026}
\externaldocument[xxii-egd-2026-]{%
xxii-egd-2026}
\externaldocument[pde-]{pde}
\externaldocument[probability-]{probability}

\externaldocument[linguistics-]{linguistics}
\externaldocument[genealogia-familiar-]{%
genealogia-familiar}

\externaldocument[%
16th-alga-meeting-2026-]{%
16th-alga-meeting-2026}

\externaldocument[%
iv-jornadas-lefschetz-2026-]{%
iv-jornadas-lefschetz-2026}

\externaldocument[reading-the-masters-]{%
reading-the-masters}

\externaldocument[real-analysis-]{%
real-analysis}

\externaldocument[optimal-transport-]{%
optimal-transport}

\externaldocument[history-]{history}

% Theorem environments.
%
\theoremstyle{plain}
\newtheorem{theorem}[subsection]{Theorem}
\newtheorem{proposition}[subsection]{Proposition}
\newtheorem{lemma}[subsection]{Lemma}

\theoremstyle{definition}
\newtheorem{definition}[subsection]{Definition}
\newtheorem{example}[subsection]{Example}
\newtheorem{exercise}[subsection]{Exercise}
\newtheorem{situation}[subsection]{Situation}

\theoremstyle{remark}
\newtheorem{remark}[subsection]{Remark}
\newtheorem{remarks}[subsection]{Remarks}

\numberwithin{equation}{subsection}

% Macros
%
\def\lim{\mathop{\mathrm{lim}}\nolimits}
\def\colim{\mathop{\mathrm{colim}}\nolimits}
\def\Spec{\mathop{\mathrm{Spec}}}
\def\Hom{\mathop{\mathrm{Hom}}\nolimits}
\def\Ext{\mathop{\mathrm{Ext}}\nolimits}
\def\SheafHom{\mathop{\mathcal{H}\!\mathit{om}}\nolimits}
\def\SheafExt{\mathop{\mathcal{E}\!\mathit{xt}}\nolimits}
\def\Sch{\mathit{Sch}}
\def\Mor{\mathop{\mathrm{Mor}}\nolimits}
\def\Ob{\mathop{\mathrm{Ob}}\nolimits}
\def\Sh{\mathop{\mathit{Sh}}\nolimits}
\def\NL{\mathop{N\!L}\nolimits}
\def\CH{\mathop{\mathrm{CH}}\nolimits}
\def\proetale{{pro\text{-}\acute{e}tale}}
\def\etale{{\acute{e}tale}}
\def\QCoh{\mathit{QCoh}}
\def\Ker{\mathop{\mathrm{Ker}}}
\def\Im{\mathop{\mathrm{Im}}}
\def\Coker{\mathop{\mathrm{Coker}}}
\def\Coim{\mathop{\mathrm{Coim}}}

% Boxtimes
%
\DeclareMathSymbol{\boxtimes}{\mathbin}{AMSa}{"02}

%
% Macros for moduli stacks/spaces
%
\def\QCohstack{\mathcal{QC}\!\mathit{oh}}
\def\Cohstack{\mathcal{C}\!\mathit{oh}}
\def\Spacesstack{\mathcal{S}\!\mathit{paces}}
\def\Quotfunctor{\mathrm{Quot}}
\def\Hilbfunctor{\mathrm{Hilb}}
\def\Curvesstack{\mathcal{C}\!\mathit{urves}}
\def\Polarizedstack{\mathcal{P}\!\mathit{olarized}}
\def\Complexesstack{\mathcal{C}\!\mathit{omplexes}}
% \Pic is the operator that assigns to X its picard group, usage \Pic(X)
% \Picardstack_{X/B} denotes the Picard stack of X over B
% \Picardfunctor_{X/B} denotes the Picard functor of X over B
\def\Pic{\mathop{\mathrm{Pic}}\nolimits}
\def\Picardstack{\mathcal{P}\!\mathit{ic}}
\def\Picardfunctor{\mathrm{Pic}}
\def\Deformationcategory{\mathcal{D}\!\mathit{ef}}

\begin{document}

\title{Homological mirror symmetry}
\maketitle
\tableofcontents

\section{Starting point}
\label{section-starting-point}

\noindent
Started on 28 September 2026 after
Sergey Galkin's message about quantum
tori and matrix factorizations.
These introductory notes were drafted by
ChatGPT. Research questions below are
not established results.

The relevant background in HP is
symplectic geometry (Poisson brackets),
geometric prequantization (commutators),
homological algebra (complexes), and
derived categories of sheaves.

\section{Where this sits}
\label{section-where-this-sits}

\noindent
Homological mirror symmetry relates
categories from symplectic geometry
to categories from algebraic geometry.
On the symplectic side, Fukaya categories
involve Lagrangian submanifolds and
Floer theory; these require further study.
On the algebraic side, one encounters
derived categories of coherent sheaves
and categories of matrix factorizations.

\begin{definition}
\label{definition-lg-model}
A \textit{Landau--Ginzburg model}, in
the setting used here, is a complex
manifold or smooth algebraic variety
$Y$ equipped with a holomorphic or
regular function
$$
W:Y\longrightarrow\mathbb{C}.
$$
The function $W$ is its
\textit{superpotential}.
\end{definition}

\noindent
A Fano mirror may be a Landau--Ginzburg
model. A del Pezzo surface is a smooth
projective surface with ample
anticanonical bundle.
There are two directions to distinguish:
coherent sheaves on the Fano side relate
to a symplectic category of its mirror;
the Fukaya category of the Fano relates
to matrix factorizations on the mirror,
with suitable choices and hypotheses.
This is an orientation, not a theorem
asserted in this generality.

Quantum cohomology modifies the cup
product using counts of rational curves.
Deformation quantization modifies
multiplication of functions into a
noncommutative product. These are
different constructions, even though
mirror symmetry relates their settings.

\section{Tori and Mumford degenerations}
\label{section-tori-and-mumford}

\begin{definition}
\label{definition-algebraic-torus}
The \textit{algebraic torus} of
dimension $n$ over $\mathbb{C}$ is
$$
T=(\mathbb{C}^{*})^n.
$$
Its coordinate algebra is
$$
\mathcal{O}(T)=
\mathbb{C}[x_1^{\pm1},\ldots,x_n^{\pm1}].
$$
\end{definition}

\noindent
In toric geometry one embeds this torus
as a dense open subset of a variety,
with an extending torus action.
Fans describe normal toric varieties.
Toric geometry supplies building blocks
for Mumford degenerations of abelian
varieties.

A quantum torus instead deforms the
multiplication of the Laurent algebra.
Thus the common starting point is the
algebraic torus, but the constructions
are different. A degeneration parameter
in Mumford's construction is not
automatically a quantization parameter.

\begin{definition}
\label{definition-quantum-torus}
For $q\in\mathbb{C}^{*}$, the
two-dimensional \textit{quantum torus}
is the algebra
$$
A_q=
\mathbb{C}\langle x^{\pm1},y^{\pm1}\rangle
/(xy-qyx).
$$
\end{definition}

\noindent
At $q=1$ this is the ordinary
commutative Laurent algebra.
For general $q$ it is an algebra,
not the coordinate ring of an ordinary
commutative affine variety.

\section{First quantum calculation}
\label{section-first-quantum-calculation}

\noindent
Work formally with $q=e^\hbar$.
The relation gives
$$
xy-yx=(e^\hbar-1)yx.
$$
Dividing by $\hbar$ and reducing
modulo $\hbar$ yields
$$
\{x,y\}=xy.
$$
Here the bracket is defined by the
first-order commutator of lifts.
This convention uses the commutator
without a factor of $i$.
Compare the Dirac rule in the
prequantization notebook, where a
different normalization is used.

\begin{exercise}
\label{exercise-laurent-bracket}
Use the Leibniz rule to show that
$$
\{x^a y^b,x^c y^d\}
=(ad-bc)x^{a+c}y^{b+d}
$$
for integers $a,b,c,d$.
\end{exercise}

\section{Matrix factorizations}
\label{section-matrix-factorizations}

\begin{definition}
\label{definition-matrix-factorization}
Let $R$ be a commutative ring and
$W\in R$. A finite free
\textit{matrix factorization} of $W$
consists of finite free modules
$E^0,E^1$ and maps
$$
d_0:E^0\longrightarrow E^1,
\qquad
d_1:E^1\longrightarrow E^0
$$
such that
$$
d_1d_0=W I_{E^0},
\qquad
d_0d_1=W I_{E^1}.
$$
Equivalently, an odd endomorphism $D$
of $E^0\oplus E^1$ satisfies
$$
D^2=W I.
$$
\end{definition}

\noindent
Compare a complex, whose differential
squares to zero. Reducing the maps
modulo $W$ gives a two-periodic complex
over $R/(W)$.
Defining the morphism complexes and the
homotopy category is a next step.

\begin{example}
\label{example-xy-factorization}
Over $R=\mathbb{C}[x,y]$, take
$$
D=\begin{pmatrix}0&x\\y&0\end{pmatrix}.
$$
Then
$$
D^2=\begin{pmatrix}xy&0\\0&xy\end{pmatrix}.
$$
Over the quantum algebra, the same
matrix multiplication gives instead
$$
D^2=\begin{pmatrix}xy&0\\0&yx\end{pmatrix}.
$$
This shows why the naive construction
does not transfer unchanged.
The commutative example is on the
affine plane, not on the torus:
on the torus $xy$ is a unit and
this factorization is contractible.
\end{example}

\noindent
For a left module over a noncommutative
algebra, multiplication by $W$ need
not be module-linear. Centrality of
$W$ guarantees linearity.
Failure of centrality obstructs the
ordinary definition; it does not prove
that every generalized version is
impossible.

\section{Galkin's question}
\label{section-galkin-question}

\noindent
On 28 September 2026, Sergey proposed:
\begin{enumerate}
\item Check whether Brown--Wemyss
normal-form methods can be applied
to elements of a quantum torus.
\item Investigate a $q$-deformation
of the mirror matrix factorization
category, provisionally denoted
$\operatorname{MF}(W_q)$.
\item If such a category can be
defined, seek its interpretation
in symplectic geometry of del Pezzo
surfaces.
\end{enumerate}

\noindent
Brown--Wemyss work with completed free
algebras and cyclic derivatives.
A quantum torus already imposes
relations and invertibility, so their
results do not apply automatically.
Their paper is a possible source of
methods, not a construction of the
category requested here.

First tasks: understand one ordinary
potential and its matrix factorizations;
choose an explicit $W_q$; check its
centrality; specify what category and
classical limit are intended.

\section{People and papers to follow}
\label{section-people-and-papers}

\noindent
Status checked on 28 September 2026.
These are reading records, not claims
that the papers have been understood.

\begin{itemize}
\item Victor: Daniel identifies his
PhD topic as matrix factorizations.
The likely full-name match is
Victor Ibrahim Santos El Adji,
Galkin's student working on graph
potentials. No matching arXiv paper
was located in this check.
This does not establish that none exists.
Confirm the exact title or arXiv link.

\item Alex Junior Gomez Saltachin:
\textit{Exceptional collections for
canonical stacks of log del Pezzo
surfaces with $\frac13(1,1)$
singularities},
\href{https://arxiv.org/abs/2606.18238}
{arXiv:2606.18238}.
The arXiv page checked lists version 4,
30 August 2026. His website uses the
title \textit{Derived categories and
Brauer groups of log del Pezzo surfaces
with 1/3(1,1) singularities} for the
same identifier; do not count these as
two different papers.

\item Alex's website also lists
\textit{Categorical genus for log del
Pezzo surfaces with cyclic quotient
singularities},
\href{https://arxiv.org/abs/2609.13102}
{arXiv:2609.13102}.
The listing was verified on his
\href{https://alequisgs.github.io/}
{personal website}; direct retrieval of
this arXiv page failed in this check.
\end{itemize}

\section{Initial reading}
\label{section-initial-reading}

\begin{itemize}
\item V. I. Arnold,
\textit{Local Normal Forms of Functions},
Inventiones mathematicae 35 (1976),
87--109.
\href{https://webhomes.maths.ed.ac.uk/%
~v1ranick/papers/arnold15.pdf}{PDF}.
Classical local singularity theory.

\item Gavin Brown and Michael Wemyss,
\textit{Local Normal Forms of
Noncommutative Functions},
\href{https://arxiv.org/abs/2111.05900}
{arXiv:2111.05900}.
Noncommutative normal forms.

\item Denis Auroux, Ludmil Katzarkov,
and Dmitri Orlov,
\textit{Mirror symmetry for Del Pezzo
surfaces: Vanishing cycles and
coherent sheaves}.
\par
\href{https://arxiv.org/abs/math/0506166}
{arXiv:math/0506166}.
A reference for the classical
del Pezzo mirror framework, not a
solution of Galkin's quantum question.

\item Denis Auroux,
\textit{Mirror symmetry and T-duality
in the complement of an anticanonical
divisor},
\href{https://arxiv.org/abs/0706.3207}
{arXiv:0706.3207}.
Geometric motivation for mirror
potentials.
\end{itemize}

\end{document}
